\documentclass[11pt]{article}

% ---------- Basic packages ----------
\usepackage[T1]{fontenc}
\usepackage{lmodern}
\usepackage{microtype}
\usepackage[margin=1in,headheight=15pt]{geometry}
\usepackage{graphicx}
\usepackage{amsmath,amssymb,amsthm}
\usepackage{mathtools}
\usepackage{tikz}
\usetikzlibrary{arrows.meta}
\usepackage{pgfplots}
\usepgfplotslibrary{groupplots}
\pgfplotsset{compat=1.18}
\usepackage[most]{tcolorbox}
\usepackage{xcolor}
\usepackage{fancyhdr}
\usepackage{titlesec}
\usepackage{enumitem}
\usepackage[hidelinks]{hyperref}
\usepackage{cleveref}
\usepackage{float}
\setlength{\parindent}{0pt}
\setlength{\parskip}{6pt}

% ---------- Document information ----------
\author{Ellie Su}
\hypersetup{
    pdftitle={Integral Inequalities},
    pdfauthor={Ellie Su}
}

% ---------- Colors ----------
\definecolor{deepblue}{HTML}{315A78}
\definecolor{softblue}{HTML}{EFF7FC}
\definecolor{purple}{HTML}{745A8D}
\definecolor{softpurple}{HTML}{F5F0FA}
\definecolor{green}{HTML}{527A63}
\definecolor{softgreen}{HTML}{EFF8F3}
\definecolor{pink}{HTML}{A76578}
\definecolor{softpink}{HTML}{FFF2F5}
\definecolor{gold}{HTML}{A78032}
\definecolor{softyellow}{HTML}{FFF9E8}
\definecolor{softgray}{HTML}{F7F7F7}
\definecolor{grayborder}{HTML}{888888}

% ---------- Theorem heading styles (no parentheses around titles) ----------
\newtheoremstyle{plaincolon}
    {6pt}{6pt}{\itshape}{}{\bfseries}{.}{0.5em}
    {\thmname{#1}\thmnumber{ #2}\thmnote{: #3}}
\newtheoremstyle{definitioncolon}
    {6pt}{6pt}{}{}{\bfseries}{.}{0.5em}
    {\thmname{#1}\thmnumber{ #2}\thmnote{: #3}}
\newtheoremstyle{remarkcolon}
    {6pt}{6pt}{}{}{\itshape}{.}{0.5em}
    {\thmname{#1}\thmnumber{ #2}\thmnote{: #3}}

% ---------- Theorem environments ----------
\theoremstyle{plaincolon}
\newtheorem{theorem}{Theorem}[section]
\newtheorem{lemma}[theorem]{Lemma}
\newtheorem{proposition}[theorem]{Proposition}
\newtheorem{corollary}[theorem]{Corollary}
\theoremstyle{definitioncolon}
\newtheorem{definition}[theorem]{Definition}
\newtheorem{example}[theorem]{Example}
\newtheorem{problem}{Problem} % independent numbering: 1, ..., 10
\newtheorem{salvage}[theorem]{Salvage}
\theoremstyle{remarkcolon}
\newtheorem{remark}[theorem]{Remark}

% ---------- Proof-style environments ----------
\newenvironment{disproof}{\begin{proof}[Disproof]}{\end{proof}}
\newenvironment{solution}{\begin{proof}[Solution]}{\end{proof}}

% ---------- General box style ----------
\tcbset{
    mathematicalbox/.style={
        enhanced,
        breakable,
        boxrule=0.8pt,
        arc=2.5mm,
        left=10pt,
        right=10pt,
        top=8pt,
        bottom=8pt,
        before skip=10pt,
        after skip=10pt,
        fonttitle=\bfseries
    }
}

% ---------- General-purpose box ----------
\newtcolorbox{mathbox}[1][]{
    mathematicalbox,
    colback=white,
    colframe=deepblue,
    coltitle=deepblue,
    borderline west={3pt}{0pt}{deepblue},
    #1
}

% ---------- Theorem and proposition boxes ----------
\tcolorboxenvironment{theorem}{
    mathematicalbox,
    colback=softblue,
    colframe=deepblue,
    borderline west={3pt}{0pt}{deepblue}
}
\tcolorboxenvironment{proposition}{
    mathematicalbox,
    colback=softblue,
    colframe=deepblue,
    borderline west={3pt}{0pt}{deepblue}
}
% ---------- Lemma boxes ----------
\tcolorboxenvironment{lemma}{
    mathematicalbox,
    colback=softpurple,
    colframe=purple,
    borderline west={3pt}{0pt}{purple}
}
% ---------- Corollary boxes ----------
\tcolorboxenvironment{corollary}{
    mathematicalbox,
    colback=softgreen,
    colframe=green,
    borderline west={3pt}{0pt}{green}
}
% ---------- Definition boxes ----------
\tcolorboxenvironment{definition}{
    mathematicalbox,
    colback=softpink,
    colframe=pink,
    borderline west={3pt}{0pt}{pink}
}
% ---------- Example boxes ----------
\tcolorboxenvironment{example}{
    mathematicalbox,
    colback=softgreen,
    colframe=green,
    borderline west={3pt}{0pt}{green}
}
% ---------- Problem boxes ----------
\tcolorboxenvironment{problem}{
    mathematicalbox,
    colback=softyellow,
    colframe=gold,
    borderline west={3pt}{0pt}{gold}
}
% ---------- Salvage boxes ----------
\tcolorboxenvironment{salvage}{
    mathematicalbox,
    colback=softpurple,
    colframe=purple,
    borderline west={3pt}{0pt}{purple}
}
% ---------- Remark boxes ----------
\tcolorboxenvironment{remark}{
    enhanced,
    breakable,
    colback=white,
    colframe=grayborder,
    boxrule=0.7pt,
    arc=2.5mm,
    left=10pt,
    right=10pt,
    top=8pt,
    bottom=8pt,
    before skip=10pt,
    after skip=10pt,
    frame style={dash pattern=on 3pt off 2pt}
}
% ---------- Proof, disproof, and solution boxes ----------
\tcolorboxenvironment{proof}{
    enhanced,
    breakable,
    colback=softgray,
    colframe=grayborder,
    boxrule=0.5pt,
    arc=2mm,
    left=10pt,
    right=10pt,
    top=8pt,
    bottom=8pt,
    before skip=8pt,
    after skip=10pt,
    borderline west={2pt}{0pt}{grayborder}
}

% ---------- Section headings ----------
\titleformat{\section}
    {\Large\bfseries\color{deepblue}}
    {\thesection.}{0.5em}{}
    [\vspace{2pt}\color{deepblue}\titlerule]
\titleformat{name=\section,numberless}
    {\Large\bfseries\color{deepblue}}
    {}{0pt}{}
    [\vspace{2pt}\color{deepblue}\titlerule]
\titleformat{\subsection}
    {\large\bfseries\color{purple}}
    {\thesubsection.}{0.5em}{}
\titleformat{name=\subsection,numberless}
    {\large\bfseries\color{purple}}
    {}{0pt}{}
\titlespacing*{\section}{0pt}{18pt}{8pt}
\titlespacing*{\subsection}{0pt}{14pt}{5pt}

% ---------- Lists ----------
\setlist[enumerate]{itemsep=3pt,topsep=5pt,leftmargin=*}
\setlist[itemize]{itemsep=3pt,topsep=5pt,leftmargin=*}

% ---------- Header and footer ----------
\pagestyle{fancy}
\fancyhf{}
\fancyhead[L]{\small\textcolor{deepblue}{Integral Inequalities}}
\fancyhead[R]{\small\textcolor{deepblue}{Ellie Su}}
\fancyfoot[C]{\small\textcolor{grayborder}{\thepage}}
\renewcommand{\headrulewidth}{0.4pt}
\renewcommand{\footrulewidth}{0pt}

% ---------- Custom title ----------
\newcommand{\settitle}[2]{%
    \thispagestyle{fancy}
    \begingroup
    \centering
    \vspace*{-0.65cm}
    {\color{deepblue}\rule{0.72\textwidth}{1.1pt}\par}
    \vspace{0.35cm}
    {\fontsize{30}{34}\selectfont\bfseries\color{deepblue}#1\par}
    \vspace{0.2cm}
    {\small\color{grayborder}#2\par}
    \vspace{0.35cm}
    {\color{deepblue}\rule{0.72\textwidth}{0.7pt}\par}
    \vspace{0.5cm}
    \endgroup
}

% ---------- Notation ----------
\numberwithin{equation}{section}
\newcommand{\CC}{\mathbb{C}}
\newcommand{\ZZ}{\mathbb{Z}}
\newcommand{\RR}{\mathbb{R}}
\newcommand{\ee}{\mathrm{e}}
\newcommand{\norm}[1]{\left\lVert#1\right\rVert}
\newcommand{\ip}[2]{\left\langle#1,#2\right\rangle}
\newcommand{\mean}[1]{\overline{#1}}
\DeclareMathOperator{\sgn}{sgn}

\begin{document}
\settitle{Integral Inequalities}{Ellie Su \quad | \quad June 2026}

\begin{mathbox}
\textbf{The central idea.} Many inequalities for finite sums continue to hold when sums are replaced with integrals. The main challenge is to recognize the right structure: a square, a product, a convex function, or a pair of similarly ordered functions. Equality cases are often as important as the inequality itself.

Unless otherwise specified, functions are real-valued and continuous on the indicated interval.
\end{mathbox}

\section{Nonnegativity and Completing the Square}

\begin{proposition}[Nonnegativity of squares]
For any continuous real-valued function $h$ on $[a,b]$,
\[
\int_a^b h(x)^2\,dx\ge 0,
\]
with equality if and only if $h(x)=0$ for every $x\in[a,b]$.
\end{proposition}

The principle becomes powerful when we can rewrite a complicated expression as a constant minus the square of another expression.

\begin{example}[A sharp quadratic bound]
For any continuous real-valued function $f$ on $[0,1]$, show that
\[
\int_0^1 f(x)(x-f(x))\,dx\le\frac1{12},
\]
and determine when equality holds.
\end{example}

\begin{solution}
Complete the square pointwise:
\[
f(x)(x-f(x))=\frac{x^2}{4}-\left(f(x)-\frac{x}{2}\right)^2.
\]
Thus
\[
\int_0^1 f(x)(x-f(x))\,dx
=\frac1{12}-\int_0^1\left(f(x)-\frac{x}{2}\right)^2dx
\le\frac1{12}.
\]
Equality holds precisely when $f(x)=x/2$ for all $x\in[0,1]$.
\end{solution}

\section{The Cauchy--Schwarz Inequality}

\begin{theorem}[Integral Cauchy--Schwarz]
For square-integrable real-valued functions $f$ and $g$ on $[a,b]$,
\[
\left(\int_a^b f(x)g(x)\,dx\right)^2\le
\left(\int_a^b f(x)^2\,dx\right)
\left(\int_a^b g(x)^2\,dx\right).
\]
Equality holds when $f$ and $g$ are proportional almost everywhere, including the case where one is zero.
\end{theorem}

\begin{proof}
For any $t\in\RR$,
\[
0\le \int_a^b(f-tg)^2dx
=\int_a^b f^2dx-2t\int_a^b fg\,dx+t^2\int_a^b g^2dx.
\]
If $\int_a^b g^2dx>0$, this quadratic in $t$ must have nonpositive discriminant, yielding the inequality. If $g=0$ almost everywhere, the result is immediate.
\end{proof}

\begin{corollary}[A fixed integral gives an $L^2$ bound]
Writing $L=b-a$ and taking $g\equiv1$,
\[
\int_a^b f(x)^2\,dx\ge\frac1L\left(\int_a^b f(x)\,dx\right)^2.
\]
Equality for continuous $f$ occurs precisely when $f$ is constant.
\end{corollary}

\begin{corollary}[Weighted Cauchy--Schwarz]
If $w(x)>0$ on $[a,b]$, then
\[
\int_a^b\frac{f(x)^2}{w(x)}\,dx\ge
\frac{\left(\int_a^b f(x)\,dx\right)^2}{\int_a^b w(x)\,dx}.
\]
Equality occurs if and only if $f(x)=c\,w(x)$ for a constant $c$.
\end{corollary}

\begin{example}[Optimizing a weighted integral]
Suppose $\int_0^1 f(x)\,dx=1$. Find the minimum of
\[
\int_0^1\frac{f(x)^2}{1+x^2}\,dx.
\]
\end{example}

\begin{solution}
Use $w(x)=1+x^2$. Then
\[
\int_0^1\frac{f(x)^2}{1+x^2}\,dx\ge
\frac{1}{\int_0^1(1+x^2)\,dx}=\frac34.
\]
Equality requires $f(x)=c(1+x^2)$. The integral constraint gives $c=3/4$, so the minimum is $\boxed{3/4}$.
\end{solution}

\section{Jensen's Inequality and Convexity}

\begin{definition}[Convexity]
A function $\phi$ on an interval is \emph{convex} if
\[
\phi(\lambda u+(1-\lambda)v)\le
\lambda\phi(u)+(1-\lambda)\phi(v)
\]
for all $u,v$ in its domain and $0\le\lambda\le1$. For twice differentiable $\phi$, the condition $\phi''\ge0$ implies convexity. A \emph{concave} function satisfies the reverse inequality.
\end{definition}

\begin{theorem}[Jensen's integral inequality]
Let $L=b-a$ and suppose $\phi$ is convex on an interval containing the range of $f$. Then
\[
\phi\left(\frac1L\int_a^b f(x)\,dx\right)
\le\frac1L\int_a^b\phi(f(x))\,dx.
\]
Reverse the inequality if $\phi$ is concave. If $\phi$ is strictly convex or strictly concave, equality for continuous $f$ requires $f$ to be constant.
\end{theorem}

\begin{remark}[Why Jensen works]
Let $m=L^{-1}\int_a^b f(x)\,dx$. For a differentiable convex $\phi$, its tangent line at $m$ lies below its graph:
\[
\phi(t)\ge\phi(m)+\phi'(m)(t-m).
\]
Set $t=f(x)$ and integrate. The average of $f(x)-m$ is zero.
\end{remark}

\begin{example}[A concave logarithm]
If $f\ge0$ and $\int_0^1 f(x)\,dx=1$, prove
\[
\int_0^1\log(1+f(x))\,dx\le\log2.
\]
\end{example}

\begin{solution}
The function $\phi(t)=\log(1+t)$ is strictly concave for $t\ge0$ because
$\phi''(t)=-(1+t)^{-2}<0$. Jensen gives
\[
\int_0^1\log(1+f(x))\,dx
\le\log\left(1+\int_0^1 f(x)\,dx\right)=\log2.
\]
Equality holds exactly for $f\equiv1$.
\end{solution}

\begin{corollary}[A useful power inequality]
For $p>1$, convexity of $|t|^p$ implies
\[
\int_a^b |f(x)|^p\,dx\ge
L^{1-p}\left|\int_a^b f(x)\,dx\right|^p.
\]
This also follows from H\"older's inequality.
\end{corollary}

\section{Chebyshev's Integral Inequality}

\begin{theorem}[Chebyshev's integral inequality]
If $f,g$ are both nondecreasing or both nonincreasing on $[a,b]$, then
\[
L\int_a^b f(x)g(x)\,dx\ge
\left(\int_a^b f(x)\,dx\right)
\left(\int_a^b g(x)\,dx\right),\qquad L=b-a.
\]
If one function is nondecreasing and the other is nonincreasing, the inequality reverses.
\end{theorem}

\begin{proof}
When $f$ and $g$ are similarly ordered,
$[f(x)-f(y)][g(x)-g(y)]\ge0$. Integrate over $[a,b]^2$ and expand:
\[
0\le\int_a^b\!\int_a^b
[f(x)-f(y)][g(x)-g(y)]\,dx\,dy
=2L\int_a^b fg-2\left(\int_a^b f\right)\left(\int_a^b g\right).
\]
Rearranging proves the claim. The opposite-monotonicity case follows by reversing the pointwise inequality.
\end{proof}

\begin{example}[Monotonicity creates a lower bound]
If $f$ is nondecreasing on $[0,1]$, prove
\[
\int_0^1 xf(x)\,dx\ge\frac12\int_0^1 f(x)\,dx.
\]
\end{example}

\begin{solution}
Apply Chebyshev to $g(x)=x$. Then
\[
\int_0^1 xf(x)\,dx\ge
\left(\int_0^1x\,dx\right)\left(\int_0^1 f(x)\,dx\right)
=\frac12\int_0^1 f(x)\,dx.
\]
Because $x$ is strictly increasing, equality for continuous nondecreasing $f$ requires $f$ to be constant.
\end{solution}

\section{How to Choose a Technique}

\begin{mathbox}
\begin{itemize}
\item \textbf{An average and a square:} try Cauchy--Schwarz with $g\equiv1$.
\item \textbf{An expression $f^2/w$:} try weighted Cauchy--Schwarz.
\item \textbf{A convex or concave function of $f$:} try Jensen.
\item \textbf{Two monotone factors:} try Chebyshev.
\item \textbf{A quadratic expression in $f$:} try completing the square.
\end{itemize}
\textbf{Always check the equality case.} It tells you both whether the estimate is sharp and what the extremizing function should be.
\end{mathbox}

\section{Practice Problems}

\emph{The problems are ordered approximately by difficulty. Unless specified otherwise, functions are continuous on $[0,1]$. Determine equality cases whenever an inequality is given.}

\subsection{Foundations}

\begin{problem}[Fixed average]
Suppose $\int_0^1 f(x)\,dx=2$. Prove
\[
\int_0^1 f(x)^2\,dx\ge4.
\]
\end{problem}

\begin{problem}[A weighted minimum]
Suppose $\int_0^1 f(x)\,dx=1$. Prove
\[
\int_0^1\frac{f(x)^2}{1+x}\,dx\ge\frac23.
\]
\end{problem}

\begin{problem}[Power integrals]
Let $a,b,t>0$. Compare
\[
P=\left(\int_0^t x^{a+b}\,dx\right)^2,
\qquad
Q=\left(\int_0^t x^{2a}\,dx\right)
\left(\int_0^t x^{2b}\,dx\right).
\]
Which is greater, and when are they equal? \emph{(Adapted from Cambridge NRICH.)}
\end{problem}

\subsection{Choosing the Inequality}

\begin{problem}[Jensen in disguise]
Assume $f(x)\ge0$ and $\int_0^1 f(x)\,dx=1$. Show
\[
\int_0^1\frac{1}{1+f(x)}\,dx\ge\frac12.
\]
\end{problem}

\begin{problem}[A reciprocal product]
Suppose $f(x)>0$ on $[0,1]$. Prove
\[
\left(\int_0^1 xf(x)\,dx\right)
\left(\int_0^1\frac{x}{f(x)}\,dx\right)\ge\frac14.
\]
\end{problem}

\begin{problem}[A monotonicity constraint]
If $f$ is nondecreasing, prove
\[
\int_0^1 x^2f(x)\,dx\ge\frac13\int_0^1 f(x)\,dx.
\]
\end{problem}

\begin{problem}[Reflection]
For a real-valued continuous function $f$, show
\[
\int_0^1 f(x)f(1-x)\,dx\le\int_0^1 f(x)^2\,dx.
\]
\end{problem}

\subsection{Challenges}

\begin{problem}[Two moment constraints]
Suppose $f$ is continuous and
\[
\int_0^1 f(x)\,dx=1,\qquad
\int_0^1 xf(x)\,dx=1.
\]
Prove $\int_0^1 f(x)^2\,dx\ge4$.
\emph{(From \emph{Putnam and Beyond}, Problem 475.)}
\end{problem}

\begin{problem}[From derivatives to functions]
Suppose $u$ is continuously differentiable and $u(0)=0$. Prove
\[
\int_0^1 u(x)^2\,dx\le\frac12\int_0^1u'(x)^2\,dx.
\]
\emph{Hint: Write $u(x)$ as an integral of its derivative.}
\end{problem}

\begin{problem}[An olympiad functional equation]
Determine all continuous $f:[0,1]\to\RR$ satisfying
\[
\int_0^1 f(x)\,dx=\frac13+\int_0^1 f(x^2)^2\,dx.
\]
\emph{(1982 Romanian Mathematical Olympiad, as presented in \emph{Putnam and Beyond}.)}
\end{problem}

\subsection*{Sources for Selected Problems}
\small
Cambridge NRICH, \emph{Integral Inequality} (Problem 3),
\href{https://nrich.maths.org/problems/integral-inequality}{NRICH online problem}.
R. Gelca and T. Andreescu, \emph{Putnam and Beyond} (Springer, 2007), integral inequalities section (Problems 8 and 10). Other problems are standard or independently formulated.
\normalsize
\end{document}